\documentclass[12pt,a4paper]{article}
\usepackage[utf8]{inputenc}
\usepackage[indonesian]{babel}
\usepackage{amsmath}
\usepackage{amsfonts}
\usepackage{amssymb}
\usepackage{geometry}

\geometry{margin=2.5cm}

\title{Menghitung Luas Permukaan Benda Putar Sumbu-y}
\author{Ahmad Fakhrul Bawani}
\date{}

\begin{document}

\maketitle

\section*{Soal}

Find the area of the surface generated by revolving the curve:
\begin{equation*}
  x = 4\sqrt{9-y}, \quad \text{untuk } 0 \le y \le \frac{143}{16}, \quad \text{mengelilingi sumbu-}y
\end{equation*}

\subsection*{1. Rumus Umum}
Rumus luas permukaan ($S$) benda putar yang dihasilkan oleh perputaran mengelilingi sumbu-$y$ untuk fungsi dalam variabel $y$ dari $y = c$ sampai $y = d$ didefinisikan sebagai:
\begin{equation*}
  S = 2\pi \int_{c}^{d} x \cdot \sqrt{1 + \left(\frac{dx}{dy}\right)^2} \, dy
\end{equation*}

\subsection*{2. Menentukan Turunan Fungsi}
Tulis ulang fungsi ke dalam bentuk pangkat: $x = 4(9-y)^{\frac{1}{2}}$. Cari turunan pertama $x$ terhadap $y$ menggunakan aturan rantai:
\begin{equation*}
  \frac{dx}{dy} = 4 \cdot \frac{1}{2}(9-y)^{-\frac{1}{2}} \cdot (-1) = \frac{-2}{\sqrt{9-y}}
\end{equation*}

Kuadratkan hasil turunan tersebut:
\begin{equation*}
  \left(\frac{dx}{dy}\right)^2 = \left(\frac{-2}{\sqrt{9-y}}\right)^2 = \frac{4}{9-y}
\end{equation*}

Masukkan ke dalam komponen akar:
\begin{equation*}
  \sqrt{1 + \left(\frac{dx}{dy}\right)^2} = \sqrt{1 + \frac{4}{9-y}} = \sqrt{\frac{9-y+4}{9-y}} = \sqrt{\frac{13-y}{9-y}} = \frac{\sqrt{13-y}}{\sqrt{9-y}}
\end{equation*}

\subsection*{3. Solusi Integrasi}
Substitusikan komponen $x = 4\sqrt{9-y}$ dan bentuk akar di atas ke dalam rumus luas permukaan dengan batas $c = 0$ dan $d = \frac{143}{16}$:
\begin{align*}
  S &= 2\pi \int_{0}^{\frac{143}{16}} \left(4\sqrt{9-y}\right) \cdot \left(\frac{\sqrt{13-y}}{\sqrt{9-y}}\right) \, dy
\end{align*}
Suku $\sqrt{9-y}$ pada pembilang dan penyebut saling menghilangkan, sehingga integralnya menjadi sangat sederhana:
\begin{align*}
  S &= 8\pi \int_{0}^{\frac{143}{16}} \sqrt{13-y} \, dy \\
  S &= 8\pi \int_{0}^{\frac{143}{16}} (13-y)^{\frac{1}{2}} \, dy
\end{align*}

Gunakan metode substitusi:
\begin{align*}
  \text{Misalkan: } & u = 13 - y \implies du = -dy \implies dy = -du \\
  \text{Batas baru: } & \text{Untuk } y = 0 \implies u = 13 - 0 = 13 \\
  & \text{Untuk } y = \frac{143}{16} \implies u = 13 - \frac{143}{16} = \frac{208 - 143}{16} = \frac{65}{16}
\end{align*}

Substitusikan variabel baru $u$ ke dalam integral:
\begin{align*}
  S &= 8\pi \int_{13}^{\frac{65}{16}} u^{\frac{1}{2}} \cdot (-du) \\
  S &= 8\pi \int_{\frac{65}{16}}^{13} u^{\frac{1}{2}} \, du \\
  S &= 8\pi \left[ \frac{2}{3}u^{\frac{3}{2}} \right]_{\frac{65}{16}}^{13} \\
  S &= \frac{16\pi}{3} \left[ u\sqrt{u} \right_{\frac{65}{16}}^{13} \\
  S &= \frac{16\pi}{3} \left( 13\sqrt{13} - \left(\frac{65}{16}\right)\sqrt{\frac{65}{16}} \right) \\
  S &= \frac{16\pi}{3} \left( 13\sqrt{13} - \frac{65\sqrt{65}}{64} \right)
\end{align*}

Kalikan $\frac{16\pi}{3}$ ke dalam kurung untuk menyederhanakan pecahan:
\begin{align*}
  S &= \frac{208\pi\sqrt{13}}{3} - \frac{16 \cdot 65\pi\sqrt{65}}{3 \cdot 64} \\
  S &= \frac{208\pi\sqrt{13}}{3} - \frac{65\pi\sqrt{65}}{12}
\end{align*}

Samakan penyebutnya menjadi 12:
\begin{equation*}
  S = \frac{\pi}{12} \left( 832\sqrt{13} - 65\sqrt{65} \right)
\end{equation*}

Jadi, luas permukaan benda putar tersebut adalah \textbf{$\dfrac{\pi}{12} \left( 832\sqrt{13} - 65\sqrt{65} \right)$ satuan luas}.

\end{document}